\documentclass[a4paper,11pt]{article}
\usepackage{exam-style-341}
\examinehead{341 农业综合知识三 · 网络技术与应用}{模拟试卷（一）}
\begin{document}

\examcover{341 农业综合知识三}{网络技术与应用 · 模拟试卷（一）}{网络技术与应用（50 分）}{50 分}{60 分钟}{本科目只考计算机网络，与程序设计、数据库无关。}

\parttitle{网络技术与应用（50 分）}

\qtypename{一、单项选择题}{10 题，每题 2 分，共 20 分}

\q 下列选项中，属于网络性能指标「时延带宽积」定义的是（\quad）
\choices{A. 链路长度与传播速度之比}{B. 发送时延与传播时延之和}{C. 传播时延与链路带宽的乘积}{D. 链路带宽与往返时间之比}

\q 在 ISO/OSI 参考模型中，负责在相邻结点之间实现无差错数据传输的是（\quad）
\choices{A. 物理层}{B. 数据链路层}{C. 网络层}{D. 运输层}

\q 按覆盖范围由小到大排列，正确的是（\quad）
\choices{A. 局域网、个人区域网、城域网、广域网}{B. 个人区域网、局域网、城域网、广域网}{C. 城域网、局域网、广域网、个人区域网}{D. 个人区域网、城域网、局域网、广域网}

\q 下列 IP 地址中属于 B 类地址的是（\quad）
\choices{A. 128.11.3.31}{B. 190.11.3.31}{C. 201.11.3.31}{D. 226.11.3.31}

\q 下列传输媒体中，抗电磁干扰能力最强、传输带宽最大的是（\quad）
\choices{A. 双绞线}{B. 同轴电缆}{C. 光纤}{D. 无线电波}

\q 关于模拟信号与数字信号，下列说法中错误的是（\quad）
\choices{A. 数字信号比模拟信号更便于加密}{B. 模拟信号在传输中失真后难以恢复}{C. 数字信号不需要任何转换即可在普通电话线上远距离传输}{D. 模拟数据可以经 PCM 转换为数字信号后再传输}

\q 下列协议中，用于把 IP 地址解析为 MAC 地址的是（\quad）
\choices{A. ARP}{B. RARP}{C. ICMP}{D. IGMP}

\q 下列关于电子邮件协议的叙述中，正确的是（\quad）
\choices{A. SMTP 用于接收邮件，POP3 用于发送邮件}{B. 用户代理把邮件交付给发送方邮件服务器使用 SMTP}{C. POP3 使用 TCP 的 25 端口}{D. IMAP 不能在服务器上管理邮件}

\q 体系结构中「对等层之间所交换的数据单元」称为（\quad）
\choices{A. 服务数据单元 SDU}{B. 协议数据单元 PDU}{C. 服务访问点 SAP}{D. 接口控制信息 ICI}

\q 关于网络层 IP 协议的特点，下列叙述正确的是（\quad）
\choices{A. 面向连接、可靠交付}{B. 无连接、尽最大努力交付、不保证可靠}{C. 提供端到端的流量控制}{D. 保证分组按序到达}

\qtypename{二、填空题}{5 题，每题 2 分，共 10 分}

\q 网络协议由三个要素组成，即语法、\fillblank{}和\fillblank{}。

\q 按覆盖范围划分，计算机网络可分为个人区域网、局域网、城域网和\fillblank{}。

\q 常用的信道复用技术有频分复用、时分复用、\fillblank{}和\fillblank{}。

\q A 类 IP 地址的第一个字节取值范围是\fillblank{}，其默认（标准）子网掩码是\fillblank{}。

\q 在 TCP 中，用来衡量报文段往返时间、并作为确定超时重传时间 RTO 主要依据的量是\fillblank{}。

\qtypename{三、名词解释}{4 题，每题 2.5 分，共 10 分}

\q 基带信号\anslines{1.2cm}

\q ARP\anslines{1.2cm}

\q FTP\anslines{1.2cm}

\q RTT\anslines{1.2cm}

\qtypename{四、简答与计算题}{2 题，每题 5 分，共 10 分}

\q 简述 CSMA/CD 协议的工作过程，并说明为什么要规定一个「争用期」以及以太网最短帧长的作用。\anslines{4cm}

\q 某公司分配到一个 B 类网络地址 \(168.195.0.0\)，现需要将其划分为 27 个子网，请回答：
\begin{enumerate}[leftmargin=2em, itemsep=2pt, topsep=3pt]
  \item 需要从主机号中借用几位作为子网号？新的子网掩码是多少？
  \item 每个子网最多可以容纳多少台主机？为什么不是 \(2^{11}=2048\) 台？
  \item 写出子网掩码的二进制形式；若不使用子网号全 0 的子网，则第 1 个可用子网的网络地址和可分配给主机的地址范围各是什么？
\end{enumerate}
\anslines{6cm}

\end{document}
